\chapter{Infrared Thermometers}

\section*{Definition and Overview}
Infrared thermometers measure body temperature without touching the skin, by detecting the infrared energy emitted by the body. The most common types are ear (tympanic) thermometers, which measure heat from the eardrum and ear canal, and forehead (temporal artery) thermometers, which scan the skin over the temporal artery. They are quick, hygienic, and easy to use, which makes them popular for children and for screening in clinics, schools, and workplaces. Accuracy depends on correct technique and the site measured, and readings may differ slightly from those taken orally or rectally. Understanding how to use and interpret them ensures reliable temperature measurement.

\section*{Epidemiology}
Infrared thermometers are among the most widely used temperature devices in homes and healthcare settings. Ear thermometers are used widely for children because they are fast and non-invasive, while forehead thermometers have become ubiquitous in screening at borders, workplaces, and events, particularly during infectious disease outbreaks. Body temperature varies with age, time of day, activity, and measurement site, which affects the interpretation of readings. Measurement errors, often from poor technique or environmental factors, are common causes of inaccurate readings and unnecessary concern.

\section*{Aetiology and Pathophysiology}
Body temperature is regulated by the hypothalamus and reflects the balance of heat produced and lost. During infection or inflammation, chemicals called pyrogens reset the set point higher, producing fever. Infrared thermometers detect the heat radiated from a surface: the eardrum shares blood supply with the hypothalamus and closely reflects core temperature, which is why ear measurement is regarded as reliable; forehead scanners measure skin temperature, which can be influenced by sweating, air movement, and ambient temperature. Because infrared devices measure surface radiation, they are sensitive to technique, position, and the environment in ways that contact thermometers are not.

\section*{Clinical Features}
\begin{itemize}
    \item Normal body temperature is around 36.1 to 37.2 degrees Celsius, varying with site and time of day
    \item Ear thermometers: a small probe placed gently in the ear canal
    \item Forehead thermometers: an infrared scan across the forehead, some touch the skin, others measure from a distance
    \item Fever is generally considered a reading of 38 degrees or above
    \item Readings may be influenced by earwax, ear infection, incorrect positioning, and cold or hot environments
    \item Different devices and sites give slightly different readings
    \item Screening uses a threshold to flag possible fever for further assessment
\end{itemize}

\section*{Differential Diagnosis}
A measured temperature must be interpreted with its limitations in mind.
\begin{itemize}
    \item \textbf{Site differences:} rectal readings are highest, then ear, oral, and skin or axillary readings
    \item \textbf{Environmental effects:} cold rooms, wind, sweating, and recent exercise alter skin readings
    \item \textbf{Ear problems:} earwax or ear infection can affect ear thermometer readings
    \item \textbf{Device error:} poor calibration, dirty lenses, or incorrect mode
    \item \textbf{Technique error:} angle, distance, or movement during measurement
    \item \textbf{True fever:} confirmed by repeat measurement or a more reliable site
\end{itemize}

\section*{When to Seek Medical Care}
A measured temperature alone is not an illness. Seek medical advice for fever that is persistent, very high, or associated with concerning symptoms, and in babies, older people, pregnant women, and people with weakened immunity. A fever in a baby under three months warrants urgent assessment. If the reading does not match the child's condition, recheck it or measure at another site.

\textbf{Call 000 (or local emergency services) for:}
\begin{itemize}
    \item A fever in a baby under three months of age
    \item A seizure or convulsion
    \item Fever with a stiff neck, severe headache, sensitivity to light, or a rash that does not fade under pressure
    \item Difficulty breathing, confusion, or a very unwell appearance
    \item Signs of dehydration: dry mouth, no wet nappies, sunken fontanelle, or lethargy
\end{itemize}

\section*{Diagnosis and Investigations}
\begin{itemize}
    \item \textbf{Device selection:} an appropriate thermometer for the age and situation
    \item \textbf{Technique:} following the manufacturer's instructions for the device
    \item \textbf{Ear thermometer use:} a snug but gentle fit in the ear canal, pulling the ear gently back
    \item \textbf{Forehead thermometer use:} correct distance and positioning as instructed
    \item \textbf{Cross-checking:} repeating the measurement or using another site when the reading is unexpected
    \item \textbf{Calibration:} regular checks according to the manufacturer's guidance
    \item \textbf{Clinical assessment:} interpreting the temperature with the person's symptoms and appearance
\end{itemize}

\section*{Management}
\begin{itemize}
    \item \textbf{Correct measurement:} follow instructions and allow the device to adjust to the room temperature
    \item \textbf{Repeated measurement:} confirm fever with more than one reading
    \item \textbf{Consider the person, not just the number:} treat the child, not the thermometer
    \item \textbf{Fever management:} fluids, light clothing, and paracetamol or ibuprofen as advised
    \item \textbf{Seek care when indicated:} for vulnerable people and concerning symptoms
    \item \textbf{Device care:} keep the sensor clean and store the device within its temperature range
    \item \textbf{Screening protocols:} follow the defined thresholds and escalation advice in workplaces and facilities
\end{itemize}

\section*{Complications}
\begin{itemize}
    \item Inaccurate readings causing missed fever and delayed care
    \item False alarms causing unnecessary worry and visits
    \item Ear injury from incorrect thermometer insertion
    \item Cross-infection if the device is not cleaned between uses
    \item Misinterpretation of skin temperature as core temperature
    \item Over-reliance on a single reading rather than clinical assessment
\end{itemize}

\section*{Prognosis and Outcomes}
Infrared thermometers are accurate and reliable when used correctly, with ear thermometers closely reflecting core temperature and forehead scanners adequate for screening. Repeated or cross-site measurement resolves most discrepancies. When readings are interpreted in the context of the person's symptoms and appearance, infrared thermometers support effective fever detection and monitoring in homes, clinics, and community settings.

\section*{Prevention and Self-Care}
\begin{itemize}
    \item Read and follow the instructions for your device
    \item Measure in a stable environment, away from heaters, fans, and direct sun
    \item Keep the ear and forehead clean and free of moisture
    \item Use the correct site for the age and situation
    \item Recheck unusual readings before acting on them
    \item Clean the device between users according to instructions
    \item Seek medical advice for fever in babies under three months and in vulnerable people
    \item Treat the child's comfort and symptoms rather than the number alone
\end{itemize}

\section*{Sources and Further Reading}
The following reputable sources provide further information for healthcare students and patients seeking reliable, current advice about thermometers and fever measurement.
\begin{itemize}
    \item Healthdirect Australia. \textit{Fever in children and how to measure temperature.} https://www.healthdirect.gov.au/
    \item Raising Children Network. \textit{Taking a child's temperature.} https://raisingchildren.net.au/
    \item Royal Children's Hospital Melbourne. \textit{Fever and temperature measurement.} https://www.rch.org.au/
    \item NHS. \textit{How to take a temperature.} https://www.nhs.uk/
    \item Mayo Clinic. \textit{Thermometers: how to take a temperature.} https://www.mayoclinic.org/
    \item Therapeutic Goods Administration. \textit{Medical thermometers regulation.} https://www.tga.gov.au/
\end{itemize}
